% Generated draft: requires mathematical and source review.
\documentclass[11pt]{article}
\usepackage[margin=25mm]{geometry}
\usepackage{fontspec}
\setmainfont{Latin Modern Roman}
\usepackage{amsmath,amssymb,mathtools,unicode-math}
\setmathfont{Latin Modern Math}
\usepackage{enumitem,needspace,xurl,hyperref}
\setlength{\parindent}{0pt}
\setlength{\parskip}{5pt}
\setlength{\emergencystretch}{3em}
\newcommand{\N}{\mathbb{N}}
\newcommand{\Z}{\mathbb{Z}}
\newcommand{\Q}{\mathbb{Q}}
\newcommand{\R}{\mathbb{R}}
\newcommand{\C}{\mathbb{C}}
\newcommand{\F}{\mathbb{F}}
\newcommand{\A}{\mathbb{A}}
\newcommand{\PP}{\mathbb{P}}
\newcommand{\E}{\mathbb{E}}
\newcommand{\eps}{\varepsilon}
\newcommand{\dd}{\,\mathrm{d}}
\begin{document}
\section*{UnsolvedMath ID: 30000379}
\textbf{Source problem code:} catalogue entry
\subsection{Definitions and mathematical statement}
Let $H$ be a Borel group in the source's category. Polishability means existence of a Polish group topology on $H$ inducing its specified Borel sets; this does not merely mean that one originally displayed topology is Polish. The excerpt does not specify whether $H$ is presented as a Borel subgroup of an ambient Polish group or as an abstract Borel group.

On $(2^{\N})^{\N}$ define
\[
 x\ E_1\ y\quad\Longleftrightarrow\quad
 \exists N\ \forall n\geq N:\ x_n=y_n.
\]
Thus $E_1$ is eventual equality of entire Cantor-space entries, not the eventual-bit relation on one copy of $2^{\N}$. For an allowed action of $H$ on a Polish space $X$, put
\[
 a\ E_H^X\ b\quad\Longleftrightarrow\quad
 \exists h\in H:\ h\cdot a=b.
\]
Borel reducibility $E_1\leq_B E_H^X$ means that there is a Borel map $F:(2^{\N})^{\N}\to X$ with
\[
 xE_1y\quad\Longleftrightarrow\quad F(x)E_H^XF(y)
 \quad\text{for all }x,y.
\]
It requires preservation and reflection of equivalence, not merely a homomorphism in one direction.

The proposed equivalence, quantified over all groups in the intended Borel-group class, is
\[
 H\text{ is not Polishable}\quad\Longleftrightarrow\quad
 \exists\text{ Polish }H\text{-space }X:\ E_1\leq_B E_H^X.
\]
The definition of Polish $H$-space is material for nonpolishable $H$: a continuous action needs a specified group topology, whereas a Borel action needs a specified product Borel structure. The excerpt does not fix that convention, and neither is silently substituted.

The supplied assessment checked 21 August 2026 attributes to Kechris and Louveau the obstruction for Polishable groups, giving the implication from existence of such a reduction to nonpolishability. It reports the converse for abelian Borel groups and for increasing unions of Polishable subgroups, but not the full equivalence. Exact action, subgroup and union hypotheses of those results are not reproduced, and their dated scope is not independently verified here.
\subsection{Short English statement}
Is a Borel group nonpolishable exactly when some permitted action on a Polish space has an orbit relation at least as complicated as eventual equality of sequences of Cantor-space points? The supplied assessment reports one direction generally and converse results for specified subclasses, while action conventions remain to be fixed.
\subsection{Sources}
\begin{itemize}
\item Local UnsolvedMath registry, permanent ID 30000379, exact source code catalogue entry; \emph{Nonpolishability and Borel Reducibility of Orbit Relations}. Proposed year 2005, status partially\_solved; corrected\_verified statement restores the group, sequence space and equivalence notation.
\item Supplied report citation: \emph{Set Theory (2006)}, Oberwolfach Reports, DOI 10.4171/owr/2005/55, pp. 3121--3174, \url{https://doi.org/10.4171/owr/2005/55}. Title year differs from DOI/proposed year.
\item Supplied literature triage checked 2026-08-21: ``original Oberwolfach report'', \url{https://doi.org/10.4171/owr/2005/55}; ``Coset equivalence relation and Polishability'', \url{https://citeseerx.ist.psu.edu/document?doi=cc534bf238cc3f905a2bc9994f369c5d10e6af21&repid=rep1&type=pdf}. It attributes the Polishable-group obstruction to Kechris and Louveau and reports the two converse subclasses. The dated OWR-derived curation requests independent verification; no linked PDF was accessed here.
\end{itemize}
\end{document}
